% =========================================================
% บท: แรงเสมือนในกรอบอ้างอิงเร่งและแรงต้านของไหล
% ต้องใช้แพ็กเกจเดิมของโครงการ เช่น
% amsmath, amssymb, graphicx, siunitx, tcolorbox
% และ environment: examplebox, exercisebox
% =========================================================

\chapter{แรงเสมือนในกรอบอ้างอิงเร่งและแรงต้านของไหล}

\begin{center}
    {\Large\bfseries Accelerated Frames, Fictitious Forces, and Fluid Resistance}

    \vspace{0.5em}

    {\large เนื้อหาเก็บตกสำหรับเตรียมสอบ สอวน. ฟิสิกส์ ค่าย 1}

    \vspace{1em}

    Tames Thanakrit Damduan
\end{center}

\vspace{1em}

% =========================================================
\section*{เป้าหมายของบทเรียน}

เมื่อเรียนจบบทนี้ นักเรียนควรสามารถ

\begin{enumerate}
    \item แยกกรอบอ้างอิงเฉื่อยและกรอบอ้างอิงเร่งได้
    \item สร้างแรงเสมือนในกรอบที่เคลื่อนที่ด้วยความเร่งได้อย่างถูกทิศ
    \item ใช้แนวคิดความเร่งโน้มถ่วงปรากฏแก้โจทย์ลิฟต์และยานพาหนะที่กำลังเร่งได้
    \item ระบุทิศของแรงต้านจากความเร็วสัมพัทธ์ระหว่างวัตถุกับของไหลได้
    \item แยกแบบจำลองแรงต้านที่แปรตาม $v$ และ $v^2$ ได้
    \item หาอัตราเร็วปลายจากเงื่อนไขแรงลัพธ์เป็นศูนย์ได้
    \item ประยุกต์แนวคิดเหล่านี้กับข้อสอบจริง สอวน. ได้
\end{enumerate}

\begin{tcolorbox}[title=ภาพรวมของบท]
บทนี้เป็นบทสั้นสำหรับรวบรวมหัวข้อที่ต่อยอดจากกฎของนิวตัน

\[
\boxed{\text{กรอบอ้างอิงเร่ง}
\quad \longrightarrow \quad
\text{เพิ่มแรงเสมือน}}
\]

\[
\boxed{\text{การเคลื่อนที่ในของไหล}
\quad \longrightarrow \quad
\text{เพิ่มแรงต้านที่ตรงข้ามกับความเร็วสัมพัทธ์}}
\]

ทั้งสองหัวข้อยังใช้หลักพื้นฐานเดิมคือ

\[
\sum \vec F=m\vec a
\]

แต่ต้องเลือกกรอบอ้างอิงและแบบจำลองแรงให้ถูกต้อง
\end{tcolorbox}

% =========================================================
\section{กรอบอ้างอิงเฉื่อยและกรอบอ้างอิงเร่ง}

\subsection{กรอบอ้างอิงเฉื่อย}

กรอบอ้างอิงเฉื่อยคือกรอบที่กฎการเคลื่อนที่ของนิวตันใช้ได้โดยตรง

\[
\boxed{\sum \vec F_{\text{real}}=m\vec a}
\]

ถ้าแรงลัพธ์เป็นศูนย์ วัตถุจะไม่มีความเร่งในกรอบนี้

\[
\sum \vec F_{\text{real}}=0
\quad \Rightarrow \quad
\vec a=0
\]

พื้นโลกสามารถถือเป็นกรอบเฉื่อยโดยประมาณสำหรับโจทย์ทั่วไปที่ไม่ต้องคำนึงถึงการหมุนของโลก

\subsection{กรอบอ้างอิงเร่ง}

รถที่เพิ่มความเร็ว รถที่เบรก ลิฟต์ที่กำลังเร่ง และยานพาหนะที่กำลังเลี้ยว เป็นตัวอย่างของกรอบอ้างอิงไม่เฉื่อย

ในบทนี้จะพิจารณากรอบที่เลื่อนที่ด้วยความเร่ง

\[
\vec A
\]

เมื่อเทียบกับกรอบเฉื่อย โดยสมมติว่าแกนพิกัดของกรอบไม่หมุน

% =========================================================
\section{ที่มาของแรงเสมือน}

ให้ตำแหน่งของวัตถุในกรอบเฉื่อยเป็น

\[
\vec r
\]

ตำแหน่งจุดกำเนิดของกรอบเร่งเป็น

\[
\vec R
\]

และตำแหน่งของวัตถุเมื่อวัดจากกรอบเร่งเป็น

\[
\vec r'
\]

จึงมีความสัมพันธ์

\[
\vec r=\vec R+\vec r'
\]

หาอนุพันธ์สองครั้งเทียบกับเวลา

\[
\vec a=\vec A+\vec a'
\]

โดย

\[
\vec A=\frac{d^2\vec R}{dt^2}
\]

คือความเร่งของกรอบ และ

\[
\vec a'=\frac{d^2\vec r'}{dt^2}
\]

คือความเร่งของวัตถุที่ผู้สังเกตในกรอบเร่งวัดได้

ในกรอบเฉื่อย กฎข้อที่สองของนิวตันให้

\[
\sum \vec F_{\text{real}}=m\vec a
\]

แทน

\[
\vec a=\vec A+\vec a'
\]

จะได้

\[
\sum \vec F_{\text{real}}
=
m(\vec A+\vec a')
\]

จัดรูป

\[
\sum \vec F_{\text{real}}-m\vec A
=
m\vec a'
\]

นิยามแรงเสมือน หรือแรงเฉื่อย เป็น

\[
\boxed{
\vec F_{\text{pseudo}}=-m\vec A
}
\]

ดังนั้น เมื่อแก้โจทย์จากกรอบเร่ง สามารถเขียนได้ว่า

\[
\boxed{
\sum \vec F_{\text{real}}
+
\vec F_{\text{pseudo}}
=
m\vec a'
}
\]

\begin{tcolorbox}[title=สมบัติของแรงเสมือน]
แรงเสมือนมีสมบัติสำคัญดังนี้

\begin{enumerate}
    \item ใช้เฉพาะเมื่อวิเคราะห์จากกรอบอ้างอิงเร่ง
    \item มีทิศตรงข้ามกับความเร่งของกรอบ
    \item มีขนาด $mA$
    \item ไม่ได้เกิดจากวัตถุอื่นมากระทำ
    \item จึงไม่มีคู่แรงตามกฎข้อที่สามของนิวตัน
\end{enumerate}

\[
\boxed{
\vec F_{\text{pseudo}}=-m\vec A
}
\]
\end{tcolorbox}

% =========================================================
\section{ความเร่งโน้มถ่วงปรากฏ}

ในกรอบที่มีความเร่ง $\vec A$ วัตถุมวล $m$ ถูกแรงโน้มถ่วง

\[
m\vec g
\]

และแรงเสมือน

\[
-m\vec A
\]

ผลรวมของสองแรงนี้เขียนได้เป็น

\[
m\vec g-m\vec A
=
m(\vec g-\vec A)
\]

จึงนิยามความเร่งโน้มถ่วงปรากฏเป็น

\[
\boxed{
\vec g_{\text{eff}}=\vec g-\vec A
}
\]

ทำให้สามารถรวมแรงโน้มถ่วงและแรงเสมือนเป็นแรงเดียวได้

\[
\boxed{
\vec F_{\text{effective gravity}}
=
m\vec g_{\text{eff}}
}
\]

\subsection{ลิฟต์เร่งขึ้น}

เลือกขึ้นเป็นทิศบวก ถ้าลิฟต์เร่งขึ้นด้วยขนาด $A$

\[
\vec A=A\hat{j}
\]

และ

\[
\vec g=-g\hat{j}
\]

ดังนั้น

\[
\vec g_{\text{eff}}
=
-g\hat{j}-A\hat{j}
\]

\[
\boxed{
\vec g_{\text{eff}}
=
-(g+A)\hat{j}
}
\]

ผู้สังเกตในลิฟต์จึงรู้สึกว่าความโน้มถ่วงมีขนาดมากขึ้นเป็น

\[
g_{\text{eff}}=g+A
\]

\subsection{ลิฟต์เร่งลง}

ถ้าลิฟต์เร่งลงด้วยขนาด $A$

\[
\boxed{
g_{\text{eff}}=g-A
}
\]

ถ้าลิฟต์ตกอิสระ

\[
A=g
\]

จะได้

\[
\boxed{
g_{\text{eff}}=0
}
\]

สภาวะนี้เรียกว่าสภาวะไร้น้ำหนักปรากฏ

% =========================================================
\begin{examplebox}{ตัวอย่าง: ลูกตุ้มในรถที่กำลังเร่ง}

ลูกตุ้มมวล $m$ แขวนอยู่ในรถที่กำลังเร่งไปทางขวาด้วยความเร่ง $A$
เมื่อมองจากรถ ลูกตุ้มอยู่นิ่งและเชือกเอียงทำมุม $\theta$
กับแนวดิ่ง จงหาความสัมพันธ์ระหว่าง $\theta$ กับ $A$

ในกรอบของรถ ต้องเพิ่มแรงเสมือนขนาด

\[
mA
\]

ชี้ไปทางซ้าย

แยกแรงตึงเชือกเป็นองค์ประกอบ

แนวราบ

\[
T\sin\theta=mA
\]

แนวดิ่ง

\[
T\cos\theta=mg
\]

หารสมการแรกด้วยสมการที่สอง

\[
\frac{T\sin\theta}{T\cos\theta}
=
\frac{mA}{mg}
\]

จึงได้

\[
\boxed{
\tan\theta=\frac{A}{g}
}
\]

เชือกเอียงไปในทิศตรงข้ามกับความเร่งของรถ
\end{examplebox}

% =========================================================
\section{แรงต้านในการเคลื่อนที่ในของไหล}

ของไหลหมายถึงสารที่สามารถไหลได้ เช่น ของเหลวและแก๊ส

เมื่อวัตถุเคลื่อนที่สัมพัทธ์กับของไหล ของไหลจะออกแรงต้านการเคลื่อนที่ของวัตถุ

ให้ความเร็วสัมพัทธ์เป็น

\[
\vec v_{\text{rel}}
=
\vec v_{\text{object}}
-
\vec v_{\text{fluid}}
\]

แรงต้านมีทิศตรงข้ามกับความเร็วสัมพัทธ์

\[
\boxed{
\vec F_D
=
-F_D(v_{\text{rel}})
\hat{u}_{\text{rel}}
}
\]

โดย $\hat{u}_{\text{rel}}$ เป็นเวกเตอร์หนึ่งหน่วยในทิศของ
$\vec v_{\text{rel}}$

\begin{tcolorbox}[title=ข้อควรระวัง]
แรงต้านไม่ได้ต้านทิศการเคลื่อนที่เมื่อเทียบกับพื้นเสมอไป
แต่ต้านการเคลื่อนที่ของวัตถุเมื่อเทียบกับของไหล

ตัวอย่างเช่น วัตถุที่อยู่นิ่งเมื่อเทียบกับพื้นอาจยังถูกแรงต้านได้
ถ้ามีกระแสลมพัดผ่าน
\end{tcolorbox}

% =========================================================
\section{แบบจำลองแรงต้านที่แปรตามอัตราเร็ว}

\subsection{แรงต้านแปรตามอัตราเร็ว}

สำหรับวัตถุขนาดเล็กที่เคลื่อนที่ช้าในของไหลหนืด
มักประมาณแรงต้านได้เป็น

\[
\boxed{
F_D=bv
}
\]

โดย $b$ เป็นค่าคงตัวที่ขึ้นกับรูปร่าง ขนาดของวัตถุ
และสมบัติของของไหล

ในรูปเวกเตอร์

\[
\boxed{
\vec F_D=-b\vec v_{\text{rel}}
}
\]

หน่วยของ $b$ หาได้จาก

\[
b=\frac{F_D}{v}
\]

ดังนั้น

\[
[b]
=
\frac{\si{N}}{\si{m.s^{-1}}}
=
\si{kg.s^{-1}}
\]

\subsection{อัตราเร็วปลายสำหรับวัตถุตก}

พิจารณาวัตถุมวล $m$ ตกในแนวดิ่ง
ให้ทิศลงเป็นบวก และให้แรงต้านมีขนาด $bv$

กฎข้อที่สองของนิวตันให้

\[
mg-bv
=
m\frac{dv}{dt}
\]

เมื่อเริ่มตก อัตราเร็วยังน้อย แรงต้านจึงน้อย
และความเร่งมีค่าใกล้ $g$

เมื่ออัตราเร็วเพิ่มขึ้น แรงต้านจะเพิ่มขึ้น
ทำให้แรงลัพธ์และความเร่งลดลง

ในที่สุดวัตถุมีอัตราเร็วคงตัว เรียกว่าอัตราเร็วปลาย

\[
v=v_T
\]

เมื่ออัตราเร็วคงตัว

\[
\frac{dv}{dt}=0
\]

ดังนั้น

\[
mg-bv_T=0
\]

จัดรูปได้

\[
\boxed{
v_T=\frac{mg}{b}
}
\]

\begin{tcolorbox}[title=หลักสำคัญของอัตราเร็วปลาย]
อัตราเร็วปลายไม่ได้หมายความว่าแรงทุกแรงหายไป

แต่หมายความว่าแรงทั้งหมดสมดุลกัน

\[
\boxed{
\sum \vec F=0
\quad \Rightarrow \quad
\vec a=0
}
\]

วัตถุจึงเคลื่อนที่ด้วยความเร็วคงตัว
\end{tcolorbox}

% =========================================================
\subsection{กรณีมีแรงลอยตัว}

วัตถุในของไหลอาจได้รับแรงลอยตัว

\[
F_B=\rho_{\text{fluid}}Vg
\]

โดย $V$ คือปริมาตรของของไหลที่ถูกแทนที่

ถ้าวัตถุกำลังลอยขึ้นด้วยอัตราเร็วปลาย
แรงลอยตัวชี้ขึ้น แต่น้ำหนักและแรงต้านชี้ลง

\[
F_B-mg-bv_T=0
\]

จึงได้

\[
\boxed{
v_T=\frac{F_B-mg}{b}
}
\]

แทน

\[
F_B=\rho_{\text{fluid}}Vg
\]

จะได้

\[
\boxed{
v_T
=
\frac{(\rho_{\text{fluid}}V-m)g}{b}
}
\]

สูตรนี้ใช้ได้เมื่อแรงลอยตัวมากกว่าน้ำหนัก

\[
\rho_{\text{fluid}}V>m
\]

% =========================================================
\section{แรงต้านแปรตามกำลังสองของอัตราเร็ว}

สำหรับวัตถุที่เคลื่อนที่เร็วในอากาศ
แรงต้านมักประมาณได้ด้วย

\[
\boxed{
F_D
=
\frac{1}{2}C_D\rho A v^2
}
\]

โดย

\[
C_D=\text{สัมประสิทธิ์แรงต้าน ซึ่งไม่มีหน่วย}
\]

\[
\rho=\text{ความหนาแน่นของของไหล}
\]

\[
A=\text{พื้นที่หน้าตัดตั้งฉากกับทิศการเคลื่อนที่}
\]

\[
v=\text{อัตราเร็วสัมพัทธ์กับของไหล}
\]

\subsection{ตรวจสอบมิติ}

พิจารณา

\[
\rho A v^2
\]

มิติของแต่ละปริมาณคือ

\[
[\rho]=ML^{-3}
\]

\[
[A]=L^2
\]

\[
[v^2]=L^2T^{-2}
\]

ดังนั้น

\[
[\rho Av^2]
=
(ML^{-3})(L^2)(L^2T^{-2})
\]

\[
[\rho Av^2]
=
MLT^{-2}
\]

ซึ่งเป็นมิติของแรง จึงสอดคล้องกัน

\subsection{อัตราเร็วปลายสำหรับการตก}

เมื่อวัตถุตกด้วยอัตราเร็วปลาย

\[
mg
=
\frac{1}{2}C_D\rho A v_T^2
\]

จัดรูป

\[
v_T^2
=
\frac{2mg}{C_D\rho A}
\]

ดังนั้น

\[
\boxed{
v_T
=
\sqrt{\frac{2mg}{C_D\rho A}}
}
\]

จากสมการนี้

\[
v_T\propto\sqrt{m}
\]

และ

\[
v_T\propto\frac{1}{\sqrt{A}}
\]

วัตถุที่มีพื้นที่หน้าตัดมากจึงมักมีอัตราเร็วปลายต่ำลง

% =========================================================
\section{การเปลี่ยนแปลงของความเร็วเมื่อแรงต้านแปรตาม $v$}

ส่วนนี้เป็นเนื้อหาอ่านเสริมสำหรับผู้ที่ต้องการเห็นรูปของ $v(t)$

สำหรับวัตถุตกจากหยุดนิ่ง ให้ทิศลงเป็นบวก

\[
m\frac{dv}{dt}=mg-bv
\]

หารด้วย $m$

\[
\frac{dv}{dt}
=
g-\frac{b}{m}v
\]

จาก

\[
v_T=\frac{mg}{b}
\]

จะเขียนได้ว่า

\[
\frac{dv}{dt}
=
\frac{b}{m}(v_T-v)
\]

แยกตัวแปร

\[
\frac{dv}{v_T-v}
=
\frac{b}{m}\,dt
\]

อินทิเกรตจาก $v=0$ ที่ $t=0$
ถึงความเร็ว $v$ ที่เวลา $t$

\[
\int_0^v
\frac{dv'}{v_T-v'}
=
\int_0^t
\frac{b}{m}\,dt'
\]

จึงได้

\[
-\ln\left(1-\frac{v}{v_T}\right)
=
\frac{bt}{m}
\]

จัดรูป

\[
1-\frac{v}{v_T}
=
e^{-bt/m}
\]

ดังนั้น

\[
\boxed{
v(t)
=
v_T\left(1-e^{-bt/m}\right)
}
\]

นิยามค่าคงตัวเวลา

\[
\boxed{
\tau=\frac{m}{b}
}
\]

จึงเขียนได้เป็น

\[
\boxed{
v(t)
=
v_T\left(1-e^{-t/\tau}\right)
}
\]

เมื่อเวลานานมาก

\[
t\rightarrow\infty
\]

จะได้

\[
e^{-t/\tau}\rightarrow0
\]

ดังนั้น

\[
v\rightarrow v_T
\]

% =========================================================
\section{วิธีเลือกแบบจำลองในการแก้โจทย์}

\begin{tcolorbox}[title=แผนแก้โจทย์กรอบเร่ง]
\begin{enumerate}
    \item ระบุว่าผู้สังเกตอยู่ในกรอบใด
    \item ถ้าเป็นกรอบเฉื่อย ใช้เฉพาะแรงจริง
    \item ถ้าเป็นกรอบเร่งด้วย $\vec A$ ให้เพิ่ม
    \[
    \vec F_{\text{pseudo}}=-m\vec A
    \]
    \item วาดแผนภาพวัตถุอิสระใหม่
    \item ตั้งสมการ
    \[
    \sum\vec F_{\text{real}}+\vec F_{\text{pseudo}}
    =m\vec a'
    \]
\end{enumerate}
\end{tcolorbox}

\begin{tcolorbox}[title=แผนแก้โจทย์แรงต้านของไหล]
\begin{enumerate}
    \item หาทิศของความเร็วสัมพัทธ์ระหว่างวัตถุกับของไหล
    \item วาดแรงต้านในทิศตรงข้ามกับความเร็วสัมพัทธ์
    \item ตรวจว่าโจทย์ให้แรงต้านแบบ $bv$ หรือ $kv^2$
    \item ใส่แรงอื่น เช่น น้ำหนักและแรงลอยตัว
    \item ถ้าเป็นอัตราเร็วปลาย ให้ใช้
    \[
    \sum\vec F=0
    \]
    แทนการใช้สูตรสำเร็จทันที
\end{enumerate}
\end{tcolorbox}

% =========================================================
\section{ชุดโจทย์รวมท้ายบท}

\begin{exercisebox}{ระดับพื้นฐาน}
\begin{enumerate}
    \item รถเร่งไปทางขวาด้วยความเร่ง $A$
    จงหาขนาดและทิศของแรงเสมือนที่กระทำต่อผู้โดยสารมวล $m$
    เมื่อวิเคราะห์จากกรอบของรถ

    \item ลิฟต์เร่งขึ้นด้วยความเร่ง $A$
    จงหาความเร่งโน้มถ่วงปรากฏที่ผู้โดยสารในลิฟต์สังเกตได้

    \item วัตถุมวล $m$ ตกในของไหลที่มีแรงต้านขนาด $bv$
    จงหาอัตราเร็วปลาย

    \item วัตถุเคลื่อนที่ไปทางตะวันออกด้วยความเร็ว $u$
    ขณะที่ลมพัดไปทางตะวันตกด้วยความเร็ว $w$
    จงหาขนาดความเร็วของวัตถุเมื่อเทียบกับอากาศ
\end{enumerate}
\end{exercisebox}

\begin{exercisebox}{ระดับกลาง}
\begin{enumerate}
    \item ลูกตุ้มแขวนในรถที่กำลังเร่งในแนวราบ
    ถ้าเชือกเอียงทำมุม $\theta$ กับแนวดิ่ง
    จงหาความเร่งของรถในรูปของ $g$ และ $\theta$

    \item วัตถุมวล $m$ และปริมาตร $V$ ลอยขึ้นในของเหลวความหนาแน่น $\rho$
    โดยมีแรงต้านขนาด $bv$
    จงหาอัตราเร็วปลาย

    \item วัตถุตกในอากาศโดยมีแรงต้าน
    \[
    F_D=\frac{1}{2}C_D\rho Av^2
    \]
    จงหาอัตราเร็วปลายในรูปของ $m$, $g$, $C_D$, $\rho$ และ $A$

    \item ร่มชูชีพทำให้พื้นที่หน้าตัดของนักกระโดดเพิ่มขึ้นเป็น $n$ เท่า
    โดยมวลและ $C_D$ คงเดิม
    จงหาอัตราส่วนของอัตราเร็วปลายใหม่ต่ออัตราเร็วปลายเดิม
\end{enumerate}
\end{exercisebox}

\begin{exercisebox}{ระดับยาก}
\begin{enumerate}
    \item กล่องมวล $m$ วางบนพื้นราบในรถที่กำลังเร่งไปทางขวาด้วยความเร่ง $A$
    สัมประสิทธิ์ความเสียดทานสถิตระหว่างกล่องกับพื้นเป็น $\mu_s$
    จงหาเงื่อนไขที่กล่องจะอยู่นิ่งเมื่อเทียบกับรถ

    \item วัตถุตกจากหยุดนิ่งในของไหลที่มีแรงต้าน $bv$
    จงพิสูจน์ว่า
    \[
    v(t)=\frac{mg}{b}
    \left(1-e^{-bt/m}\right)
    \]

    \item วัตถุทรงกลมสองลูกมีมวลเท่ากันและมี
    $C_D$ เท่ากัน แต่พื้นที่หน้าตัดของลูกที่สองเป็น
    $4$ เท่าของลูกแรก จงหาอัตราส่วนของอัตราเร็วปลาย
    \[
    \frac{v_{T,2}}{v_{T,1}}
    \]
\end{enumerate}
\end{exercisebox}

% =========================================================
\section{สรุปสูตรสำคัญ}

\begin{tcolorbox}[title=สรุปกรอบอ้างอิงเร่ง]
ความสัมพันธ์ของความเร่ง

\[
\vec a=\vec A+\vec a'
\]

แรงเสมือน

\[
\boxed{
\vec F_{\text{pseudo}}=-m\vec A
}
\]

กฎข้อที่สองเมื่อมองจากกรอบเร่ง

\[
\boxed{
\sum\vec F_{\text{real}}
+
\vec F_{\text{pseudo}}
=
m\vec a'
}
\]

ความเร่งโน้มถ่วงปรากฏ

\[
\boxed{
\vec g_{\text{eff}}=\vec g-\vec A
}
\]

ลูกตุ้มในรถที่เร่งแนวราบ

\[
\boxed{
\tan\theta=\frac{A}{g}
}
\]
\end{tcolorbox}

\begin{tcolorbox}[title=สรุปแรงต้านของไหล]
ความเร็วสัมพัทธ์

\[
\vec v_{\text{rel}}
=
\vec v_{\text{object}}
-
\vec v_{\text{fluid}}
\]

แรงต้านชี้ตรงข้ามกับ $\vec v_{\text{rel}}$

แรงต้านที่แปรตามอัตราเร็ว

\[
\boxed{
\vec F_D=-b\vec v_{\text{rel}}
}
\]

แรงต้านที่แปรตามกำลังสองของอัตราเร็ว

\[
\boxed{
F_D=\frac{1}{2}C_D\rho Av^2
}
\]

อัตราเร็วปลายเมื่อ $F_D=bv$

\[
\boxed{
v_T=\frac{mg}{b}
}
\]

อัตราเร็วปลายเมื่อ
$F_D=\frac12C_D\rho Av^2$

\[
\boxed{
v_T
=
\sqrt{\frac{2mg}{C_D\rho A}}
}
\]

กรณีวัตถุลอยขึ้นและมีแรงลอยตัว $F_B$

\[
\boxed{
v_T=\frac{F_B-mg}{b}
}
\]

การตกจากหยุดนิ่งเมื่อแรงต้านเป็น $bv$

\[
\boxed{
v(t)=v_T\left(1-e^{-t/\tau}\right)
}
\]

\[
\boxed{
\tau=\frac{m}{b}
}
\]
\end{tcolorbox}

% =========================================================
\newpage
\section{ข้อสอบจริง สอวน. ที่เกี่ยวข้องกับบทนี้}


% =========================================================
\subsection{ปี 2560 ข้อ 3: หลอดไฟตกในลิฟต์ที่เร่งขึ้น}

\vspace{1em}

\IfFileExists{img/posn60q3.png}{
\begin{center}
    \includegraphics[width=1\linewidth]{img/posn60q3.png}
\end{center}
}{
\begin{tcolorbox}[title=ตำแหน่งภาพที่แนะนำ]
ให้ตัดภาพข้อสอบปี 2560 ข้อ 3 แล้วบันทึกเป็น

\[
\texttt{img/posn60q3.png}
\]
\end{tcolorbox}
}

\begin{examplebox}{เฉลยละเอียด}

ลิฟต์เร่งขึ้นด้วยความเร่ง

\[
A=2.0\,\si{m.s^{-2}}
\]

หลอดไฟอยู่สูงจากพื้นลิฟต์

\[
h=2.95\,\si{m}
\]

เมื่อหลอดไฟหลุด หลอดไฟมีความเร็วต้นเท่ากับพื้นลิฟต์
ดังนั้นความเร็วต้นเมื่อวัดจากลิฟต์เป็นศูนย์

\[
v'_0=0
\]

พิจารณาจากกรอบของลิฟต์

แรงจริงที่กระทำต่อหลอดไฟคือแรงโน้มถ่วง

\[
mg
\]

ชี้ลง

เพราะลิฟต์เร่งขึ้น จึงต้องเพิ่มแรงเสมือน

\[
mA
\]

ชี้ลง

ดังนั้นแรงลัพธ์ในกรอบลิฟต์คือ

\[
mg+mA
\]

ความเร่งสัมพัทธ์ของหลอดไฟจึงมีขนาด

\[
a'=g+A
\]

แทนค่า

\[
a'=9.8+2.0
\]

\[
a'=11.8\,\si{m.s^{-2}}
\]

ใช้สมการการเคลื่อนที่สัมพัทธ์

\[
h=v'_0t+\frac{1}{2}a't^2
\]

เพราะ $v'_0=0$

\[
2.95
=
\frac{1}{2}(11.8)t^2
\]

\[
t^2
=
\frac{5.90}{11.8}
\]

\[
t^2=\frac{1}{2}
\]

ดังนั้น

\[
\boxed{
t=\frac{1}{\sqrt{2}}\,\si{s}
}
\]
\end{examplebox}

% =========================================================
\subsection{ปี 2560 ตอนที่ 2 ข้อ 2:
มวลติดอยู่ข้างรถเข็นที่กำลังเร่ง}

\vspace{1em}

\IfFileExists{img/posn60p2q2.png}{
\begin{center}
    \includegraphics[width=1\linewidth]{img/posn60p2q2.png}
\end{center}
}{
\begin{tcolorbox}[title=ตำแหน่งภาพที่แนะนำ]
ให้ตัดภาพข้อสอบปี 2560 ตอนที่ 2 ข้อ 2 แล้วบันทึกเป็น

\[
\texttt{img/posn60p2q2.png}
\]
\end{tcolorbox}
}

\begin{examplebox}{เฉลยละเอียด}

มวลเล็ก $m$ สัมผัสกับผนังแนวดิ่งของมวลใหญ่ $M$
โดยสัมประสิทธิ์ความเสียดทานสถิตเป็น

\[
\mu_s=0.50
\]

ให้ระบบเร่งไปทางขวาด้วยความเร่ง $a$

พิจารณาจากกรอบของมวลใหญ่ $M$
มวล $m$ อยู่นิ่งเมื่อเทียบกับกรอบนี้
จึงต้องเพิ่มแรงเสมือน

\[
ma
\]

ชี้ไปทางซ้าย

ในแนวราบ แรงปกติ $N$ ชี้ไปทางขวา
และสมดุลกับแรงเสมือน

\[
N-ma=0
\]

ดังนั้น

\[
N=ma
\]

ในแนวดิ่ง แรงเสียดทานสถิตต้องพยุงน้ำหนัก

\[
f_s=mg
\]

แต่แรงเสียดทานสถิตมีค่าสูงสุด

\[
f_{s,\max}=\mu_sN
\]

เพื่อไม่ให้มวลไถลลง ต้องมี

\[
mg\leq\mu_sN
\]

แทน

\[
N=ma
\]

จะได้

\[
mg\leq\mu_sma
\]

ตัด $m$

\[
g\leq\mu_sa
\]

ดังนั้น

\[
a\geq\frac{g}{\mu_s}
\]

ค่าความเร่งต่ำสุดคือ

\[
a_{\min}
=
\frac{g}{0.50}
\]

\[
a_{\min}=2g
\]

แรงภายนอก $F$ ต้องเร่งมวลรวม $M+m$

\[
F=(M+m)a
\]

ดังนั้น

\[
\boxed{
F_{\min}=2(M+m)g
}
\]
\end{examplebox}

% =========================================================
\subsection{ปี 2560 ข้อ 1:
มิติของค่าคงตัวในสมการแรงต้าน}

\vspace{1em}

\IfFileExists{img/posn60q1.png}{
\begin{center}
    \includegraphics[width=1\linewidth]{img/posn60q1.png}
\end{center}
}{
\begin{tcolorbox}[title=ตำแหน่งภาพที่แนะนำ]
ให้ตัดภาพข้อสอบปี 2560 ข้อ 1 แล้วบันทึกเป็น

\[
\texttt{img/posn60q1.png}
\]
\end{tcolorbox}
}

\begin{examplebox}{เฉลยละเอียด}

โจทย์ให้แรงต้านมีขนาด

\[
F=kv^2A
\]

ต้องการหาว่า $k$ มีมิติของปริมาณใด

จัดรูป

\[
k=\frac{F}{v^2A}
\]

มิติของแรงคือ

\[
[F]=MLT^{-2}
\]

มิติของอัตราเร็วยกกำลังสองคือ

\[
[v^2]=L^2T^{-2}
\]

มิติของพื้นที่คือ

\[
[A]=L^2
\]

ดังนั้น

\[
[k]
=
\frac{MLT^{-2}}
{(L^2T^{-2})(L^2)}
\]

\[
[k]
=
ML^{-3}
\]

มิตินี้ตรงกับความหนาแน่น

\[
[\rho]=ML^{-3}
\]

ดังนั้น

\[
\boxed{k\text{ มีมิติของความหนาแน่น}}
\]

\[
\boxed{\text{ตอบ A}}
\]
\end{examplebox}

% =========================================================
\subsection{ปี 2560 ข้อ 4:
อัตราเร็วปลายของลูกปิงปองที่ลอยขึ้นในน้ำ}

\vspace{1em}

\IfFileExists{img/posn60q4.png}{
\begin{center}
    \includegraphics[width=1\linewidth]{img/posn60q4.png}
\end{center}
}{
\begin{tcolorbox}[title=ตำแหน่งภาพที่แนะนำ]
ให้ตัดภาพข้อสอบปี 2560 ข้อ 4 แล้วบันทึกเป็น

\[
\texttt{img/posn60q4.png}
\]
\end{tcolorbox}
}

\begin{examplebox}{เฉลยละเอียด}

ลูกปิงปองมีมวล $m$ และปริมาตร $V$
อยู่ในน้ำความหนาแน่น $\rho$

เมื่อเชือกขาด ลูกปิงปองเคลื่อนที่ขึ้น

แรงที่กระทำมีสามแรง

แรงลอยตัวชี้ขึ้น

\[
F_B=\rho Vg
\]

น้ำหนักชี้ลง

\[
mg
\]

แรงต้านชี้ลง เพราะลูกปิงปองกำลังเคลื่อนที่ขึ้น

\[
F_D=Cv
\]

เมื่อถึงอัตราเร็วปลาย

\[
a=0
\]

ดังนั้นแรงลัพธ์เป็นศูนย์

\[
F_B-mg-Cv_T=0
\]

แทนแรงลอยตัว

\[
\rho Vg-mg-Cv_T=0
\]

จัดรูป

\[
Cv_T
=
\rho Vg-mg
\]

\[
Cv_T
=
(\rho V-m)g
\]

ดังนั้น

\[
\boxed{
v_T
=
\frac{(\rho V-m)g}{C}
}
\]

\[
\boxed{\text{ตอบ D}}
\]
\end{examplebox}
